\documentclass[11pt]{article}
\usepackage[margin=1in]{geometry}
\usepackage{amsmath, amssymb, amsthm, mathtools}
\usepackage{hyperref}

\newtheorem{theorem}{Theorem}
\newtheorem{corollary}[theorem]{Corollary}
\theoremstyle{definition}
\newtheorem{definition}[theorem]{Definition}
\newtheorem{example}[theorem]{Example}

\title{Events and Sigma-Algebras}
\author{math-foundations / concept 18}
\date{}

\begin{document}
\maketitle

\subsection*{First, an example}

Before the abstract axioms, fix $\Omega = \{1, 2, 3, 4\}$ and look at three sigma-algebras on it. The \emph{trivial} sigma-algebra is $\mathcal{F}_{0} = \{\emptyset, \Omega\}$, with just $2$ elements: it can only say ``something happened'' or ``nothing happened.'' The \emph{discrete} sigma-algebra is the full power set $2^{\Omega}$, with $2^{4} = 16$ elements: every subset is an event. In between, the sigma-algebra generated by the singleton $\{1\}$ is
\[
\sigma(\{\{1\}\}) = \{\emptyset,\ \{1\},\ \{2,3,4\},\ \Omega\},
\]
\emph{Read aloud:} sigma of the singleton 1 equals the set containing the empty set, the singleton 1, the set 2 3 4, and Omega.
with $4$ elements: it can resolve ``did $1$ occur?'' and nothing finer. These three illustrate the spectrum from coarse to fine that the formal definition below captures.

\section{Motivation}

For a finite or countable sample space $\Omega$ we may assign probabilities to every subset. For uncountable $\Omega$ this fails: by the Vitali construction there exist subsets of $[0,1]$ that admit no translation-invariant, countably-additive probability. To rescue the theory, we restrict attention to a distinguished family of \emph{measurable} subsets called a sigma-algebra.

\section{Definitions}

\begin{definition}[Sigma-algebra]
Let $\Omega$ be a non-empty set. A collection $\mathcal{F} \subseteq 2^{\Omega}$ is a \emph{sigma-algebra} on $\Omega$ if
\begin{enumerate}
    \item[(S1)] $\Omega \in \mathcal{F}$;
    \item[(S2)] for every $A \in \mathcal{F}$, the complement $A^{c} = \Omega \setminus A$ also lies in $\mathcal{F}$;
    \item[(S3)] for every sequence $A_{1}, A_{2}, \ldots \in \mathcal{F}$, the countable union $\bigcup_{n=1}^{\infty} A_{n}$ also lies in $\mathcal{F}$.
\end{enumerate}
A pair $(\Omega, \mathcal{F})$ with $\mathcal{F}$ a sigma-algebra on $\Omega$ is a \emph{measurable space}, and the elements of $\mathcal{F}$ are called \emph{events} (or \emph{measurable sets}).
\end{definition}

\begin{definition}[Generated sigma-algebra]
Let $\mathcal{C} \subseteq 2^{\Omega}$. The \emph{sigma-algebra generated by $\mathcal{C}$}, written $\sigma(\mathcal{C})$, is the smallest sigma-algebra on $\Omega$ that contains $\mathcal{C}$, i.e.
\[
\sigma(\mathcal{C}) \;=\; \bigcap \bigl\{\, \mathcal{F} \subseteq 2^{\Omega} \;:\; \mathcal{F} \text{ is a sigma-algebra on } \Omega,\ \mathcal{C} \subseteq \mathcal{F} \,\bigr\}.
\]
\emph{Read aloud:} sigma of script C equals the intersection over all script F, subset of the power set of Omega, such that script F is a sigma-algebra on Omega and script C is contained in script F.
The intersection is well-defined because $2^{\Omega}$ is itself a sigma-algebra containing $\mathcal{C}$, and because of Theorem~\ref{thm:intersection} below.
\end{definition}

\begin{definition}[Borel sigma-algebra]
The \emph{Borel sigma-algebra} on $\mathbb{R}$ is $\mathcal{B}(\mathbb{R}) := \sigma(\mathcal{T})$, where $\mathcal{T}$ is the family of open subsets of $\mathbb{R}$. Equivalently, $\mathcal{B}(\mathbb{R}) = \sigma(\{(-\infty, a] : a \in \mathbb{R}\})$.
\end{definition}

\section{Key theorem}

\begin{theorem}[Arbitrary intersection of sigma-algebras]\label{thm:intersection}
Let $\Omega$ be a non-empty set and let $\{\mathcal{F}_{\alpha}\}_{\alpha \in I}$ be a (possibly uncountable) family of sigma-algebras on $\Omega$. Then
\[
\mathcal{G} \;:=\; \bigcap_{\alpha \in I} \mathcal{F}_{\alpha}
\]
\emph{Read aloud:} script G is defined as the intersection over alpha in I of script F sub alpha.
is a sigma-algebra on $\Omega$.
\end{theorem}

\begin{proof}
We verify each of the three axioms (S1), (S2), (S3) for $\mathcal{G}$.

\medskip
\noindent\textbf{(S1) Whole space.} For every $\alpha \in I$, axiom (S1) applied to $\mathcal{F}_{\alpha}$ gives $\Omega \in \mathcal{F}_{\alpha}$. Hence $\Omega$ lies in every member of the family, so
\[
\Omega \in \bigcap_{\alpha \in I} \mathcal{F}_{\alpha} = \mathcal{G}.
\]
\emph{Read aloud:} Omega is in the intersection over alpha in I of script F sub alpha, which equals script G.

\medskip
\noindent\textbf{(S2) Closure under complement.} Let $A \in \mathcal{G}$. By definition of intersection, $A \in \mathcal{F}_{\alpha}$ for every $\alpha \in I$. Applying axiom (S2) inside each $\mathcal{F}_{\alpha}$, we get $A^{c} \in \mathcal{F}_{\alpha}$ for every $\alpha$. Therefore
\[
A^{c} \in \bigcap_{\alpha \in I} \mathcal{F}_{\alpha} = \mathcal{G}.
\]
\emph{Read aloud:} A complement is in the intersection over alpha in I of script F sub alpha, which equals script G.

\medskip
\noindent\textbf{(S3) Closure under countable union.} Let $A_{1}, A_{2}, \ldots$ be a sequence in $\mathcal{G}$. Fix any $\alpha \in I$. Since each $A_{n} \in \mathcal{G} \subseteq \mathcal{F}_{\alpha}$, the sequence $\{A_{n}\}_{n \ge 1}$ lies entirely in $\mathcal{F}_{\alpha}$. By axiom (S3) applied to $\mathcal{F}_{\alpha}$,
\[
\bigcup_{n=1}^{\infty} A_{n} \in \mathcal{F}_{\alpha}.
\]
\emph{Read aloud:} the union from n equals 1 to infinity of A sub n is in script F sub alpha.
This holds for every $\alpha \in I$, so $\bigcup_{n=1}^{\infty} A_{n} \in \bigcap_{\alpha} \mathcal{F}_{\alpha} = \mathcal{G}$.

\medskip
All three axioms hold; therefore $\mathcal{G}$ is a sigma-algebra on $\Omega$.
\end{proof}

\begin{corollary}
For every class $\mathcal{C} \subseteq 2^{\Omega}$, the generated sigma-algebra $\sigma(\mathcal{C})$ exists, is unique, and equals the intersection of all sigma-algebras on $\Omega$ that contain $\mathcal{C}$.
\end{corollary}

\begin{proof}
The family $\mathcal{S} := \{\mathcal{F} \subseteq 2^{\Omega} : \mathcal{F} \text{ is a sigma-algebra},\ \mathcal{C} \subseteq \mathcal{F}\}$ is non-empty since $2^{\Omega} \in \mathcal{S}$. By Theorem~\ref{thm:intersection}, $\bigcap_{\mathcal{F} \in \mathcal{S}} \mathcal{F}$ is a sigma-algebra. It contains $\mathcal{C}$ by construction, and is the smallest such by minimality of intersection.
\end{proof}

\section{Examples}

\begin{example}[Trivial sigma-algebra]
$\mathcal{F}_{0} = \{\emptyset, \Omega\}$ is a sigma-algebra: $\Omega \in \mathcal{F}_{0}$; $\Omega^{c} = \emptyset \in \mathcal{F}_{0}$; any countable union of $\emptyset$'s and $\Omega$'s is either $\emptyset$ or $\Omega$.
\end{example}

\begin{example}[Sigma-algebra generated by one event]
For $A \subseteq \Omega$ with $A \notin \{\emptyset, \Omega\}$,
\[
\sigma(\{A\}) = \{\emptyset, A, A^{c}, \Omega\}.
\]
\emph{Read aloud:} sigma of the singleton A equals the set containing the empty set, A, A complement, and Omega.
\end{example}

\end{document}
